\section{实现与机制}
\FocusSlide[IMPLEMENTATION]{把算法中的一次操作，\\落实为线程中的一次访问。}

\begin{frame}[fragile]{代码与机制：一个线程负责一个有效元素}
\begin{columns}[T,onlytextwidth]
\begin{column}{.58\textwidth}
\subhead{CUDA：element-wise vector addition}
\begin{lstlisting}[basicstyle=\small\ttfamily,aboveskip=2pt,belowskip=2pt]
__global__ void add(
    const float* a, const float* b,
    float* c, size_t n) {
  const size_t i =
      size_t(blockIdx.x) * blockDim.x
      + threadIdx.x;
  if (i < n) {
    c[i] = a[i] + b[i];
  }
}
\end{lstlisting}
\captiontext{代码片段用于解释索引与访问关系。调用者负责分配有效设备缓冲区并检查执行状态。}
\end{column}
\begin{column}{.38\textwidth}
\subhead{三个阅读锚点}
\textbf{索引}\par $i=\mathrm{blockIdx.x}\times B+\mathrm{threadIdx.x}$。
\par\vspace{0.2mm}\textbf{边界}\par $i<n$ 控制最后一个 block 的有效线程。
\par\vspace{0.2mm}\textbf{访问}\par 相邻线程处理连续元素，形成相邻地址访问。
\notebox{越界保护限定有效索引；输入输出缓冲区需满足契约。}
\end{column}
\end{columns}
\end{frame}

\begin{frame}{访存布局：用地址表解释连续访问}
\subhead{同一个 Warp 的前四个 Lane，32-bit 元素}
\begin{tabularx}{\textwidth}{@{}p{3.4cm}XXXX@{}}
\toprule 地址偏移（Byte） & Lane 0 & Lane 1 & Lane 2 & Lane 3 \\\midrule
Stride = 1 element & 0 & 4 & 8 & 12 \\
Stride = 32 elements & 0 & 128 & 256 & 384 \\
\bottomrule
\end{tabularx}
\vspace{5mm}
\begin{columns}[T,onlytextwidth]
\begin{column}{.48\textwidth}
\subhead{计算规则}
\[\mathrm{offset}(\ell)=4s\ell\]
\smalltext $\ell$ 为 Lane 编号，$s$ 为元素间距。地址表由公式直接计算。
\end{column}
\begin{column}{.48\textwidth}
\subhead{与实测连接}
使用内存事务、有效带宽与 Kernel 时间评估访问模式，并同时记录对齐条件。
\notebox{内存事务数量还取决于参与 Lane、访问宽度、地址对齐和 GPU 架构。}
\end{column}
\end{columns}
\end{frame}

\begin{frame}{机制说明：从局部计算走向完整执行路径}
\begin{columns}[T,onlytextwidth]
\begin{column}{.47\textwidth}
\subhead{局部关系}
\begin{denseitems}
\item 连续地址访问影响访存事务。
\item 线程状态影响 Register 需求。
\item 协作范围影响同步方式。
\end{denseitems}
\vspace{3mm}\insight{局部机制用于解释观察结果。}
\end{column}
\begin{column}{.47\textwidth}
\subhead{完整路径}
\begin{denseitems}
\item 调用与数据准备进入计时边界。
\item 传输与 Kernel 的依赖需要明确。
\item 输出完成后才能计算端到端延迟。
\end{denseitems}
\vspace{3mm}\insight{最终指标覆盖用户实际等待的过程。}
\end{column}
\end{columns}
\vspace{3mm}\notebox{当局部收益与总时间变化不一致时，使用阶段计时检查成本转移，并保持输入与配置一致。}
\end{frame}
